%===============================================================================
\section{Related Work}\label{sec:related}
%===============================================================================

\paragraph{Inference and aggregation in data release}
That individually innocuous data items can jointly reveal protected information is an old
observation in database security, where it is known as the inference and aggregation
problem~\cite{lunt1989aggregation,farkas2002inference}; the privacy literature supplies the
best-known illustration, three demographic attributes that jointly single out most
individuals~\cite{sweeney2000simple}. In power systems, public market data have been shown to
expose non-public physical information, for example grid topology from locational marginal
prices~\cite{kekatos2014topology}. Privacy mechanisms for metering data~\cite{mcdaniel2009smartgrid,giaconi2018smartmeter}
and formal models such as differential privacy~\cite{dwork2014dp} or
$k$-anonymity~\cite{sweeney2002kanonymity} bound what can be learned about an individual
record after a release mechanism is applied. They answer a different question from the one
faced by a market operator, who must decide which aggregate operational \emph{fields} may be
published at all.

\paragraph{Data sensitivity assessment under inference leakage}
Data classification and grading is performed field by field in practice, and the Chinese
national standard explicitly asks that aggregation and processing effects be
considered~\cite{gbt43697}, without prescribing how. The closest work to ours is IGNN~\cite{hu2026ignn},
which introduced data inference leakage into sensitivity assessment for cyber-physical power
grids: fields are nodes of an inference graph whose edges encode physical or statistical
inference relations, and sensitivity is propagated along these edges by a graph attention
network from an intrinsic (expert-assigned) score. This established that the risk of a field
depends on its inference relations to sensitive data. Two questions remain open in this line
of work. First, the risk score is calibrated against expert labels, while its ultimate
validation is an attack experiment in which a model is trained on a group of fields; the
attack outcome itself is not the definition. Second, an edge connects two fields, whereas
inference frequently requires several fields jointly, with none of them informative alone.
We take the attack outcome on a field \emph{set} as the definition of risk and derive the
field-level quantity from it.

\paragraph{Feature importance}
Turning set-level values into per-feature scores is the subject of feature attribution.
Shapley-based measures~\cite{shapley1953value,lundberg2017shap,covert2020sage,williamson2020spvim}
average a feature's marginal contribution over all contexts, which is the right answer for
fair allocation but dilutes the credit of redundant features~\cite{kumar2020shapley,catav2021mci}.
Leave-one-covariate-out (LOCO)~\cite{lei2018loco} and permutation importance~\cite{breiman2001rf,fisher2019mcr,hooker2021unrestricted}
evaluate a single context, the full feature set, so a feature with a near-substitute receives
almost no importance. The Marginal Contribution Index (MCI)~\cite{catav2021mci} takes the
\emph{maximum} marginal contribution over contexts, is characterized axiomatically, and is
immune to dilution. Security reasoning favours the worst case for the same reason
quantitative information flow replaced Shannon entropy by min-entropy and gain
functions~\cite{smith2009qif,alvim2012gleakage}: an average over contexts is not a guarantee.
Our field-level risk is a budget-constrained MCI evaluated relative to already-public data.

\paragraph{Fast estimation of many refitted models}
Computing the performance of a model refitted on many feature subsets is expensive, which has
motivated lazy~\cite{gao2022lazyvi} and warm-started~\cite{sun2025warmstart} estimators of
variable importance, surrogate models trained under random masks for removal-based
explanations~\cite{covert2021explaining,jethani2022fastshap,covert2023vitshapley}, conditional generative models trained under
random masks~\cite{ivanov2019vaeac,li2020acflow,yoon2020vime,he2022mae}, and tabular foundation models that predict from any
set of columns in context~\cite{hollmann2025tabpfn}.
Closed-form ridge heads on frozen features are the standard fast-adaptation device in
meta-learning~\cite{bertinetto2019r2d2,lee2019metaoptnet}, and the benefit of sharing a
representation across tasks is well understood~\cite{caruana1997multitask,baxter2000inductive,maurer2016benefit}.
We combine these ingredients into an estimator of set-level inferability and, more
importantly, use it inside a scan--certify loop so that decisions rest on retrained
attackers rather than on the estimate.

%===============================================================================
\section{Problem Formulation}\label{sec:problem}
%===============================================================================

\subsection{Parties, field roles and threat model}
The \textbf{auditor} is the data owner. It holds the complete history of all operational
fields, including those it regards as sensitive, and grades fields offline before release.
The \textbf{adversary} obtains published fields and is assumed to hold historical auxiliary
labels of the sensitive target, from which it trains a supervised inference model; risk is
therefore measured as out-of-sample \emph{supervised} inferability. The adversary may combine
the field under review with other fields it already holds.

Field roles are fixed by disclosure rules, not by data analysis. Let $G$ be the candidate
fields that the rules allow to be published, $B\subseteq G$ a \emph{public baseline} that is
already released ($B=\emptyset$ unless stated otherwise), $H=G\setminus B$ the fields under
review, and $Y_y$ a sensitive target. Section~\ref{sec:data} instantiates these roles for
four data sets from the categories and release timings of the Chinese electricity-market
disclosure rules~\cite{nea2024disclosure}. An ex-ante estimate of the \emph{same physical
quantity} as a target (e.g.\ a load forecast when the target is metered load) is removed from
$G$, because it is a copy of the target rather than an inference channel.

\subsection{Set-level inferability}
\begin{definition}[Population inferability]
\[
\Vy^{*}(S) = 1 - \inf_{f}\,\mathbb{E}\big[(Y_y - f(X_S))^{2}\big] \big/ \mathrm{Var}(Y_y),
\]
the infimum taken over measurable $f$. It satisfies $\Vy^{*}(\emptyset)=0$, $0 \le \Vy^{*} \le 1$
and is monotone in $S$.
\end{definition}

\begin{definition}[Empirical inferability]\label{def:vret}
The data are split into training and test rows, with a validation part inside training. For
an attacker family $\mathcal{H}$, every $h\in\mathcal{H}$ is trained on $X_S$ only; the
attacker and its training epoch are selected on validation, and the test $\Rtwo$ is reported:
$\Vy^{\mathrm{ret}}(S) = \Rtwo_{\mathrm{test}}(\hat h_S)$ with
$\hat h_S = \argmax_{h \in \mathcal{H},\,\mathrm{epoch}} \Rtwo_{\mathrm{val}}(h;S)$.
\end{definition}

Finite samples, optimization and selection noise make $\Vy^{\mathrm{ret}}$ slightly
non-monotone. An adversary holding $S$ can train on any subset of $S$ and keep the best, so
the operational object is the \emph{monotone closure}
\begin{equation}\label{eq:closure}
\Vy(S) = \Vy^{\mathrm{ret}}\big(\hat T_S\big), \qquad
\hat T_S = \argmax_{T \subseteq S} \Rtwo_{\mathrm{val}}(T),
\end{equation}
where the subset is chosen on validation and the value is read on test. We write $\Vy$ for
the closure throughout. It is the \emph{ground truth} of this paper: every conclusion about
risk is stated in terms of $\Vy$ computed by dedicated retraining.

\subsection{Audit questions}
The auditor asks three questions: \emph{(Q1)} if a set $S$ is published, how accurately can
$Y_y$ be inferred? \emph{(Q2)} how risky is a single field $i$, why, and together with which
other fields? \emph{(Q3)} how can Q1--Q2 be answered for all fields and targets, repeatedly as
the public baseline or target definitions change, without training a dedicated attacker for
every set? Section~\ref{sec:mk} answers Q2 in terms of $\Vy$, and Section~\ref{sec:model}
answers Q3.

%===============================================================================
\section{Field-Level Inference Risk}\label{sec:mk}
%===============================================================================

\subsection{Contextual gain}
The quantity from which everything else follows is the gain that releasing field $i$ brings over a
given background $T$ of fields the adversary already holds,
\begin{equation}\label{eq:delta}
\Delta_i(A)=\Vy(A \cup \{i\}) - \Vy(A),\qquad A=B\cup T .
\end{equation}
$\Delta_i(B)$ is the single-field risk. The same field can be nearly useless over one background
and decisive over another: a price series says little about a unit's output alone, but together
with the unit's forecast availability it reveals the residual quantity the unit covers. A
field-level risk quantity must therefore summarize a whole \emph{distribution} of contextual gains
$\{\Delta_i(B\cup T)\}_T$, and the question is which summary serves security.

\subsection{Budgeted marginal risk and background amplification}
\begin{definition}[Budgeted marginal risk, background amplification]\label{def:mk}
For target $y$, public baseline $B$, fields under review $H=G\setminus B$ and budget $K$,
\begin{equation}\label{eq:mk}
\Mik \;=\; \max_{T \subseteq H \setminus \{i\},\, |T| \le K}\ \Delta_i(B\cup T),
\qquad
\Gamma^{(K)}_{i\to y} \;=\; \Mik - M^{(0)}_{i\to y} \;\ge 0 .
\end{equation}
Any maximizer $T^{*}$ is a \emph{witness background}.
\end{definition}
$\Mik$ is the largest increase of inferability that releasing $i$ can cause when the adversary
holds at most $K$ other fields; $\Gamma^{(K)}_{i\to y}$ is the part of it that exists only because of
the background, i.e.\ the amplification of single-field risk by combination. For each (field,
target) pair the audit reports the profile $\big(M^{(0)}_{i\to y},\ \Gamma^{(K)}_{i\to y},\ \Mik\big)$
together with the top-$r$ backgrounds ranked by $\Delta_i$; in our data the top backgrounds are
often nearly tied and share a core partner field, so the partner is more informative than a
single argmax. At $K=1$ the amplification has a closed form.

\begin{proposition}[Amplification at $K=1$]\label{prop:k1}
With $I^B_{ij}=\Vy(B\cup\{i,j\})-\Vy(B\cup\{i\})-\Vy(B\cup\{j\})+\Vy(B)$,
$\Gamma^{(1)}_{i\to y}=\max\big(0,\max_{j} I^{B}_{ij}\big)$.
\end{proposition}
\begin{proof}
Over the background $\{j\}$, $\Delta_i(B\cup\{j\})=M^{(0)}_{i\to y}+I^B_{ij}$; over $\emptyset$ it is
$M^{(0)}_{i\to y}$. Take the maximum.
\end{proof}
A field is amplified by one background field exactly when it is strictly super-additive with it,
and the amplification equals the strongest positive interaction.

\subsection{Why the maximum}
Turning a set function into a per-field score is the problem of feature attribution, and the
natural first choice is an average of $\Delta_i$ over backgrounds, as in Shapley values or
SAGE~\cite{covert2020sage}. The average answers ``how much does $i$ contribute typically'', which is
the right question for explanation and fair allocation. A security audit asks instead whether
\emph{some} admissible background makes $i$ dangerous; a field whose gain is small over most
backgrounds but large over one is dangerous, and averaging hides it---in our data the average
contextual gain is typically about a third of the maximum (Section~\ref{sec:res-amp}). This is the
same reason quantitative information flow measures leakage by worst-case
vulnerability rather than by average entropy~\cite{smith2009qif,alvim2012gleakage}. The maximum
marginal contribution is the Marginal Contribution Index (MCI)~\cite{catav2021mci}; $\Mik$ adds the
two conditions a disclosure audit needs, a public baseline $B$ and a bounded background $K$.
Table~\ref{tab:req} summarizes the requirements and which field scores meet them.

\begin{table}[!tbp]
\centering\small
\caption{Requirements for a field-level inference-risk quantity. \checkmark: satisfied;
$\circ$: partially; --: not satisfied.}\label{tab:req}
\resizebox{\textwidth}{!}{%
\begin{tabular}{lcccccc}
\toprule
Method & checkable by attack & joint inference & worst case & no dilution & bounded $K$ & public baseline $B$\\
\midrule
Correlation / mutual information & -- & -- & -- & \checkmark & -- & -- \\
Single-field retraining $M^{(0)}$ & \checkmark & -- & -- & \checkmark & -- & $\circ$ \\
Inference-graph propagation~\cite{hu2026ignn} & -- & $\circ$ & -- & $\circ$ & -- & -- \\
Shapley / SAGE~\cite{covert2020sage} & \checkmark & \checkmark & -- & -- & -- & -- \\
LOCO / permutation & \checkmark & \checkmark & -- & -- & -- & -- \\
MCI~\cite{catav2021mci} & \checkmark & \checkmark & \checkmark & \checkmark & -- & -- \\
$\Mik$ (this paper) & \checkmark & \checkmark & \checkmark & \checkmark & \checkmark & \checkmark \\
\bottomrule
\end{tabular}}
\end{table}

\begin{proposition}[Relation to MCI]\label{prop:mci}
Let $v_B(T) := \Vy(B \cup T) - \Vy(B)$ for $T \subseteq H$. Then $v_B$ is monotone,
non-negative and $v_B(\emptyset)=0$; $\Mik = \max_{|T|\le K}[v_B(T\cup\{i\}) - v_B(T)]$ is the
budget-constrained MCI of $v_B$; for $K \ge |H|-1$ and $B=\emptyset$ it is the MCI of $\Vy$.
\end{proposition}
The axiomatic properties of MCI for monotone evaluation functions~\cite{catav2021mci}---dummy,
symmetry, the self-contribution lower bound $\Mik\ge M^{(0)}_{i\to y}$ and invariance to
duplicated features---therefore transfer and are cited rather than claimed; $\Mik$ is
non-decreasing in $K$ and reaches the MCI at $K=|H|-1$.

\subsection{Security meaning}
\begin{theorem}[Threshold-free safety margin]\label{thm:margin}
Fix $i$, $y$, $B$, $K$. A number $m\ge 0$ satisfies ``for every threshold $\tau$ and every
background $|T|\le K$: $\Vy(B\cup T)\le\tau-m \Rightarrow \Vy(B\cup T\cup\{i\})\le\tau$''
if and only if $m\ge\Mik$. For a fixed $\tau$, the smallest $m$ with this property can be
strictly smaller than $\Mik$.
\end{theorem}
\begin{proof}
If $m\ge\Mik$ then $\Vy(B\cup T\cup\{i\})=\Vy(B\cup T)+\Delta_i(B\cup T)\le(\tau-m)+m$.
Conversely, for any $T$ choose $\tau=\Vy(B\cup T)+m$; the property gives $\Delta_i(B\cup T)\le
m$. Strictness: with $B=\emptyset$, $K=1$, $\Vy(\{i\})=0.1$, $\Vy(\{j\})=0.5$,
$\Vy(\{i,j\})=0.9$ one has $M^{(1)}_{i\to y}=0.4$, yet at $\tau=0.95$ the necessary margin is $0$.
\end{proof}
Thus $\Mik$ is the answer to ``by how much can publishing $i$ raise the adversary's accuracy,
whatever it already holds within budget''---a property of the field that does not depend on
the threshold later chosen for grading.

\begin{definition}[Criticality and minimal unsafe sets]
Under baseline $B$, $E\subseteq H$ is a $\tau$-\emph{minimal unsafe set} (MUS) if
$\Vy(B\cup E)>\tau$ and $\Vy(B\cup E')\le\tau$ for every $E'\subsetneq E$. Field $i$ is
$K$-\emph{background} $\tau$-\emph{critical} if some $T$ with $|T|\le K$, $i\notin T$, has
$\Vy(B\cup T)\le\tau<\Vy(B\cup T\cup\{i\})$.
\end{definition}

\begin{theorem}[Criticality $\iff$ small MUS]\label{thm:mus}
If $\Vy$ is monotone and $\Vy(B)\le\tau$, field $i$ is $K$-background $\tau$-critical if and
only if it belongs to a MUS of size at most $K+1$.
\end{theorem}
\begin{proof}
($\Leftarrow$) For a MUS $E\ni i$ with $|E|\le K+1$ take $T=E\setminus\{i\}$.
($\Rightarrow$) If $\Vy(B\cup T)\le\tau<\Vy(B\cup T\cup\{i\})$, take an inclusion-minimal unsafe
$E\subseteq T\cup\{i\}$; it is a MUS of size $\le K+1$. If $i\notin E$ then $E\subseteq T$ and
monotonicity gives $\Vy(B\cup T)\ge\Vy(B\cup E)>\tau$, a contradiction.
\end{proof}
Theorem~\ref{thm:mus} makes the field-level flag a faithful projection of the set-level
danger structure: flagging the $\tau$-critical fields misses no field that takes part in a
dangerous combination of at most $K+1$ fields, and the witness $T^*$ names the combination.
The same list of minimal unsafe sets also supports downstream disposal
(Section~\ref{sec:application}).
Monotonicity is essential: without the closure, a field can cross the threshold over a
background that already contains an unsafe subset, and no MUS then contains it.


\subsection{Choice of $K$}
Different $K$ describe adversaries with different background knowledge, and $\Mik$ grows with $K$.
We report $K=2$ as the main setting and measure $K=3$ on every data set. The threat model is an
adversary holding a handful of other published series rather than the whole catalogue;
high-order interactions are rare in engineered systems~\cite{li2006regularities}, and models with
pairwise interactions are close to unrestricted ones on tabular data~\cite{lou2013ga2m}. There is,
however, no general saturation theorem---a parity function of $K+2$ fields defeats any fixed
$K$---and saturation of the maximal gain does not imply that no larger set crosses a threshold.
Section~\ref{sec:res-robust} measures both.

%===============================================================================
\section{An Amortized Attacker for Set-Level Inferability}\label{sec:model}
%===============================================================================

\subsection{Why an estimator is needed}
Evaluating $\Mik$ for all fields at $K=2$ requires $\Vy$ on every set of size at most three;
with 35 candidates this is 7,175 sets per target, each requiring a family of attackers to be
trained. The table has to be recomputed whenever the public baseline grows, a target is
redefined, or a what-if release is examined. We therefore build one model that answers
$\Vy(S)$ for any $S$ quickly (Fig.~\ref{fig:model}), and we make sure that no conclusion depends on its accuracy
alone (Section~\ref{sec:scan}).

\begin{figure}[!tbp]
\centering
\includegraphics[width=\textwidth]{fig3_model.pdf}
\caption{The amortized attacker. (a) One backbone is pre-trained under random visibility masks on the
auditor's data, to predict the targets and reconstruct the hidden columns, and frozen. (b) A query $(S,y)$ solves a ridge
readout in closed form on the features of $S$---from the pre-trained backbone and an untrained network of the same
architecture---per seed, and averages the predictions. (c) The scan proposes backgrounds per field; retraining only these
certifies a lower bound on $\Mik$.}
\label{fig:model}
\end{figure}

\subsection{Masked pre-training}
The adversary does not know in advance which fields it will obtain. We model this by
pre-training one network under random visibility masks. A mask $m\in\{0,1\}^{p}$ is drawn by
first sampling the number of visible fields uniformly and then a uniform subset of that size;
the network $f_\theta(x\odot m,\ m)$ receives the masked values and the mask itself, so that a
hidden field is distinguishable from a zero value. The objective combines predicting the sensitive targets from the visible fields with reconstructing every hidden column
from the visible ones~\cite{yoon2020vime,he2022mae},
\begin{equation}\label{eq:pretrain}
\mathcal L(\theta)=\mathbb{E}_{m}\Big[\big\|f^{\,y}_\theta(x\odot m,m)-y\big\|^2
+\gamma\,\frac{\big\|(1-m)\odot\big(f^{\,x}_\theta(x\odot m,m)-x\big)\big\|^2}{\|1-m\|_1}\Big],
\end{equation}
with $\gamma=1$. The reconstruction term is essential. With $\gamma=0$ (target-only pre-training) the network collapses its
hidden layer onto the few directions that predict the targets on average over masks, and on the small sets of an audit it
then adds nothing over an untrained network (Section~\ref{sec:res-est}); reconstruction forces the representation to encode
the operating state that drives all fields. Reconstructing all columns the auditor holds also decouples the backbone from
any particular target definition. The last hidden layer (256 units) is kept as a frozen feature map
$\phi_\theta(x\odot m, m)$; early stopping uses the target loss on validation rows only. Alongside each pre-trained backbone
we keep an untrained network of the same architecture, $\phi_{\vartheta}$ with random initial weights $\vartheta$. The two
play different roles in the readout: the pre-trained features supply target-aligned directions and the operating state, the
random features the high-dimensional bulk of a kernel regression~\cite{rahimi2007random} (Section~\ref{sec:res-est}).

\subsection{Subset-specific closed-form readout}
For a query $(S,y)$ with mask $m_S$ and seed $r$, the feature block is
\begin{equation}
Z_{r}(S) = \big[\, x\odot m_S,\ (x \odot m_S)^{2},\ \phi_{\theta_r}(x\odot m_S,\ m_S),\ \phi_{\vartheta_r}(x\odot m_S,\ m_S) \,\big],
\end{equation}
standardized with training statistics. A ridge readout
\begin{equation}\label{eq:ridge}
\beta_{r}(S) = \big(Z_r^{\top} Z_r + \lambda I\big)^{-1} Z_r^{\top} y
\end{equation}
is solved in closed form on the training rows, with $\lambda$ chosen by five-fold
cross-validation inside the training rows. Each of $R=3$ seeds (a pre-trained and an untrained network) has its own
readout---three separate closed-form solves per query, not one solve on the features of all seeds---and their
predictions are averaged, $\hat y = R^{-1}\sum_r Z_r^{\mathrm{test}}\beta_r(S)$, and
$\hat V_y(S)=\Rtwo_{\mathrm{test}}(\hat y)$, followed by the monotone closure. Numerical
safeguards matter on real data: a floor of $10^{-3}$ on feature standard deviations, and
clipping test features to their training range. Without them, units that are almost constant
on the training rows (inactive ReLUs) are amplified on test rows, a few predictions explode and
$\hat V$ collapses to zero for those sets; since $\Mik$ is a maximum over differences, such
collapses turn directly into inflated risk (Section~\ref{sec:res-est}). A reconstruction backbone
needs a third safeguard, clipping predictions to the range of the target on the training rows: an adversary never
predicts a value it has not seen. No test label enters
pre-training, the readout, the choice of $\lambda$ or the closure.

\subsection{Why the readout is solved per subset}
A cheaper alternative lets the network output the prediction directly, $\hat y=f_\theta(x\odot
m_S, m_S)$, and reads $\hat V$ off it. Then one set of output weights serves every
conditioning pattern, i.e.\ the attacker is forced to combine every subset of fields in the
same way---the very effect that combination risk is about. Solving the readout per subset
keeps the expensive representation shared and the combination rule specific.

\subsection{Error propagation}
Write the estimate as $\hat V(S)=V(S)+b+\varepsilon(S)$, where $b$ is a bias common to the sets
of interest and $\varepsilon(S)$ the set-specific part. Since a marginal is a difference of two
nested sets, $\hat\Delta_i(T)-\Delta_i(T)=\varepsilon(T\cup\{i\})-\varepsilon(T)$ and
\begin{equation}\label{eq:prop}
\big|\hat M^{(K)}_{i\to y} - \Mik\big| \;\le\; \max_{|T|\le K}\big|\varepsilon(T\cup\{i\})-\varepsilon(T)\big|,
\end{equation}
independently of $b$. A readout on a frozen representation typically underestimates the fully
trained attacker by a similar amount on nested sets, which is exactly a common-mode $b$; this is
why an estimator can be biased in $V$ yet accurate in $M$. The quantity to control is the
\emph{marginal} error, which we report throughout.

\subsection{Scan and certify}\label{sec:scan}
The estimator is used to \emph{scan}: it produces $\hat M^{(K)}$ for every field together with a
ranked list of candidate backgrounds. Decisions are then \emph{certified}: for each field,
dedicated attackers are retrained on the top-$k$ backgrounds only, and
$L^{(K)}_{i\to y}=\max_{T\in\text{top-}k}\Delta_i(B\cup T)\le\Mik$ is a lower bound computed
entirely from retrained attackers. The ratio $\sum_i L/\sum_i M$ measures how much of the true
risk the scan recovers at a certification cost of $k$ retrained backgrounds per field instead
of all $\binom{|H|-1}{\le K}$.
